%% annex/p09/csl-corrections.tex --- endpoint-sensitive CSL calculations.

\section{CSL interval endpoints and transient analysis}
\label{sec:annex-p09-csl-corrections}

\begin{annexitem}{annex:p9.csl.1}{thm:p09-csl-endpoints}{CSL until endpoint normalization}
Use the chain, state sets \(A,B\), interval \(I\), closed counterpart \(J\),
and probability function \(u_I\) from Theorem~\ref{thm:p09-csl-endpoints}.
\begin{proof}
A Poisson clock of rate \(\lambda>0\), with
\(\lambda\geq E(s)\) for every state, realizes the chain by accepting a tick
at \(s\) with probability \(E(s)/\lambda\), then choosing the actual
successor by \(P(s,\cdot)\).  Rejected ticks are not path steps; when
\(E(s)=0\), every tick is rejected.  While the state remains \(s\), the
probability that no tick has been accepted by time \(t\) is
\[
 \sum_{n\geq0}e^{-\lambda t}\frac{(\lambda t)^n}{n!}
       (1-E(s)/\lambda)^n=e^{-E(s)t},
\]
which gives the stated exponential residence law.  There are almost surely
only finitely many ticks in every bounded time interval, and no tick at a
prescribed positive time.  Therefore
paths are time-divergent and are locally constant on both sides of each
prescribed positive endpoint, outside a null event.

Work outside the null event for the positive finite endpoints of \(I\).
If an until witness occurs only at an excluded upper endpoint \(b>a\),
local constancy allows a slightly earlier witness inside \(I\).
If it occurs at an excluded lower endpoint \(a>0\), the prior-time condition
and local constancy imply both \(\Phi\) and \(\Psi\) there and for a short
time afterwards, so a slightly later witness lies in \(I\).
If \(a=0\) and \(s\in A\), the same later-witness argument works because
the initial residence time is strictly positive, or infinite at an absorbing
state.  If \(a=0\), \(0\notin I\), and \(s\notin A\), every possible
witness is positive and fails the prior-time condition at \(u=0\).
This proves the displayed cases.  A singleton has only its stated witness;
at zero the prior-time condition is empty.  An empty interval has no witness.
\end{proof}
\end{annexitem}

\begin{annexitem}{annex:p9.csl.2}{thm:p09-csl-interval-computations}{Endpoint-aware CSL until computations}
Use the fixed state sets and modified generators from
Theorem~\ref{thm:p09-csl-interval-computations}.  The no-jump-at-a-fixed-time
fact is established in Annex~\ref{annex:p9.csl.1}.
The identification of \(h\) with unbounded-until probability uses the cited
embedded-chain result in Theorem~\ref{thm:p09-csl-model-checking}.  The proof
below supplies the interval decomposition; it does not reprove the source's
Volterra or steady-state results.
\begin{proof}[Derivation of the interval computations]
The chain with generator \(Q_{A\setminus B}\) follows the original chain
until it reaches \(B\) or leaves \(A\cup B\), then remains at that success
or failure state.  With right-continuous occupation, the first entry into
\(B\) is itself an admissible until witness, provided all preceding states
belonged to \(A\).  Consequently its probability of being in \(B\) at time
\(b\) is exactly \(u_{[0,b]}\), including the immediate witness for a
starting state in \(B\).

When \(a>0\), the original path must remain in \(A\) throughout
\([0,a)\).  A jump exactly at the deterministic time \(a\) has probability
zero, so the same condition can require membership in \(A\) at \(a\).
The \(s\)-th row of \(e^{Q_Aa}D_A\) gives precisely the subprobability
distribution at time \(a\) of paths from \(s\) that have met this condition.
In this first phase, states in \(A\cap B\) are not made absorbing: a visit to
\(B\) before \(a\) is not yet a witness.  The Markov property at time
\(a\), followed by the preceding success/failure construction for duration
\(b-a\), gives the second formula.  For \(a=b>0\), it reduces to
\(e^{Q_Aa}\mathbf 1_{A\cap B}\); it is not a first-entry probability at
one prescribed instant.  For \(a=b=0\), the first formula gives
\(v_0=\mathbf 1_B\).  The same decomposition with the embedded-chain
reachability vector \(h\) proves the infinite-upper-bound formulas.
All zero-rate rows remain absorbing, so these arguments include them.
\end{proof}

\paragraph{The separate timed-next calculation.}
For the timed-next formula printed after
Theorem~\ref{thm:p09-csl-interval-computations}, the probability is zero
when \(E(s)=0\), since there is no next path step.  For \(E(s)>0\),
the factor \(e^{-E(s)a}-e^{-E(s)b}\) is the waiting-time probability,
and \(\sum_{t\in B}P(s,t)\) is the probability of an eligible successor.
This follows by
conditioning on the first successor and integrating the exponential density.
\end{annexitem}
